% GENERATED FILE. Edit canonical lessons, then run npm run generate:books.
% canonical-source-hash: fnv1a64:b25455b6e3dd1582
% canonical-lessons: HI-C237-caukor, HI-C237-tikona, HI-C237-sankara, HI-C237-uthla, HI-C237-basi

\chapter{Square, Triangular, Narrow, Shallow, Stale}
\label{ch:hi-square-triangular-narrow-shallow-stale}
\hlchaptermodality{237}

\begin{chapteropening}
\textbf{By the end of this chapter:} \emph{I can say square, triangular, narrow, shallow and stale.}

Complete the last lesson of chapter 237: I can say square, triangular, narrow, shallow and stale.
\end{chapteropening}

\section[\texorpdfstring{\hi{चौकोर} (caukor)}{चौकोर (caukor)}]{\hi{चौकोर} (caukor) — square}

\label{lesson:HI-C237-caukor}



\begin{quote}
Before the new one: say the Hindi for separate, then the Hindi for round.
\end{quote}

\subsection*{You'll want to know: \hi{चौकोर}}
\textbf{\hi{चौकोर}} — \emph{caukor} — \textquotedblleft{}square\textquotedblright{}.

\begin{grammarlens}[title={What you've built}]
One more word for this chapter.
\end{grammarlens}

\subsection*{Your turn}
\begin{itemize}
\raggedright
  \item \emph{Say it:} \emph{caukor}
  \item \emph{Say it:} \emph{caukor}, once more
  \item \emph{Say it:} say \emph{caukor}
  \item \emph{Recall:} say the Hindi for separate, then the Hindi for round, then say \emph{caukor} again
\end{itemize}

\begin{tcolorbox}[breakable,colback=teal!4,colframe=teal!35!black,title={Before you move on}]
What does \hi{चौकोर} mean? (\textquotedblleft{}Square\textquotedblright{}.) Say it once more.
\end{tcolorbox}
\section[\texorpdfstring{\hi{तिकोना} (tikoṇā)}{तिकोना (tikoṇā)}]{\hi{तिकोना} (tikoṇā) — triangular}

\label{lesson:HI-C237-tikona}



\begin{quote}
Before the new one: say the Hindi for round, then the Hindi for square.
\end{quote}

\subsection*{You'll want to know: \hi{तिकोना}}
\textbf{\hi{तिकोना}} — \emph{tikoṇā} — \textquotedblleft{}triangular\textquotedblright{}.

\begin{grammarlens}[title={What you've built}]
One more word for this chapter.
\end{grammarlens}

\subsection*{Your turn}
\begin{itemize}
\raggedright
  \item \emph{Say it:} \emph{tikoṇā}
  \item \emph{Say it:} \emph{tikoṇā}, once more
  \item \emph{Say it:} say \emph{tikoṇā}
  \item \emph{Recall:} say the Hindi for round, then the Hindi for square, then say \emph{tikoṇā} again
\end{itemize}

\begin{tcolorbox}[breakable,colback=teal!4,colframe=teal!35!black,title={Before you move on}]
What does \hi{तिकोना} mean? (\textquotedblleft{}Triangular\textquotedblright{}.) Say it once more.
\end{tcolorbox}
\section[\texorpdfstring{\hi{संकरा} (saṅkarā)}{संकरा (saṅkarā)}]{\hi{संकरा} (saṅkarā) — narrow}

\label{lesson:HI-C237-sankara}



\begin{quote}
Before the new one: say the Hindi for square, then the Hindi for triangular.
\end{quote}

\subsection*{You'll want to know: \hi{संकरा}}
\textbf{\hi{संकरा}} — \emph{saṅkarā} — \textquotedblleft{}narrow\textquotedblright{}.

\begin{grammarlens}[title={What you've built}]
One more word for this chapter.
\end{grammarlens}

\subsection*{Your turn}
\begin{itemize}
\raggedright
  \item \emph{Say it:} \emph{saṅkarā}
  \item \emph{Say it:} \emph{saṅkarā}, once more
  \item \emph{Say it:} say \emph{saṅkarā}
  \item \emph{Recall:} say the Hindi for square, then the Hindi for triangular, then say \emph{saṅkarā} again
\end{itemize}

\begin{tcolorbox}[breakable,colback=teal!4,colframe=teal!35!black,title={Before you move on}]
What does \hi{संकरा} mean? (\textquotedblleft{}Narrow\textquotedblright{}.) Say it once more.
\end{tcolorbox}
\section[\texorpdfstring{\hi{उथला} (uthlā)}{उथला (uthlā)}]{\hi{उथला} (uthlā) — shallow}

\label{lesson:HI-C237-uthla}



\begin{quote}
Before the new one: say the Hindi for triangular, then the Hindi for narrow.
\end{quote}

\subsection*{You'll want to know: \hi{उथला}}
\textbf{\hi{उथला}} — \emph{uthlā} — \textquotedblleft{}shallow\textquotedblright{}.

\begin{grammarlens}[title={What you've built}]
One more word for this chapter.
\end{grammarlens}

\subsection*{Your turn}
\begin{itemize}
\raggedright
  \item \emph{Say it:} \emph{uthlā}
  \item \emph{Say it:} \emph{uthlā}, once more
  \item \emph{Say it:} say \emph{uthlā}
  \item \emph{Recall:} say the Hindi for triangular, then the Hindi for narrow, then say \emph{uthlā} again
\end{itemize}

\begin{tcolorbox}[breakable,colback=teal!4,colframe=teal!35!black,title={Before you move on}]
What does \hi{उथला} mean? (\textquotedblleft{}Shallow\textquotedblright{}.) Say it once more.
\end{tcolorbox}
\section[\texorpdfstring{\hi{बासी} (bāsī)}{बासी (bāsī)}]{\hi{बासी} (bāsī) — stale}

\label{lesson:HI-C237-basi}



\begin{quote}
Before the new one: say the Hindi for narrow, then the Hindi for shallow.
\end{quote}

\subsection*{You'll want to know: \hi{बासी}}
\textbf{\hi{बासी}} — \emph{bāsī} — \textquotedblleft{}stale\textquotedblright{}.

\begin{grammarlens}[title={What you've built}]
One more word for this chapter.
\end{grammarlens}

\subsection*{Your turn}
\begin{itemize}
\raggedright
  \item \emph{Say it:} \emph{bāsī}
  \item \emph{Say it:} \emph{bāsī}, once more
  \item \emph{Say it:} say \emph{bāsī}
  \item \emph{Recall:} say the Hindi for narrow, then the Hindi for shallow, then say \emph{bāsī} again
\end{itemize}

\begin{tcolorbox}[breakable,colback=teal!4,colframe=teal!35!black,title={Before you move on}]
What does \hi{बासी} mean? (\textquotedblleft{}Stale\textquotedblright{}.) Say it once more.
\end{tcolorbox}
